\clearpage
\scanpage{4}
\mspage{\unsure{\handwritten{(32}|(32}}

% One long pencil cancellation descends from the right of the opening lines
% to the left of the last paragraph. Local obliterations are nested below.
\newsavebox{\coteFourTypescript}
\begin{lrbox}{\coteFourTypescript}
\begin{minipage}[t]{.96\linewidth}
\llap{\handwritten{\tikz[baseline]{\draw[thin] (0,0) .. controls (.16,.2) and (.34,.37) .. (.43,.17);}}}Considérant la restriction de $\handwritten{\varphi}$ à la sous-catégorie pleine de $\overset{\handwritten{\sim}}{\mathrm{X}_{\text{ét}}}$ formée des ouverts de ce topos, ou ce qui revient au même \handwritten{(6.1.)} des ouverts $\mathrm U$ de $\mathrm X$, on trouve une application
\[
\handwritten{\varphi_{0}}:\handwritten{\mathcal U}(\mathrm X)\mathrel{\mbox{-}\mkern-2mu\mbox{-}\mkern-3mu\handwritten{\to}}\handwritten{\bigl\{}0,1\handwritten{\bigr\}}
\]
commutant aux sup quelconques et aux inf finis. (On note que pour $\mathrm U\handwritten{\in\mathcal U}(\mathrm X)$, $\handwritten{\varphi}(\mathrm U)$ est vide ou réduit un point, et on prend $\handwritten{\varphi_{0}}(\mathrm U)=0$ ou $1$ suivant qu'on est dans l'un ou l'autre cas).

On voit facilement \deleted{(\illegible{1}\unsure{)}}, grâce au fait que tout fermé irréductible de $\mathrm X$ a un point génerique et un seul, que $\handwritten{\varphi_{0}}$ est defini à l'aide d'un \deleted{\unsure{e|x}} unique $\mathrm x\handwritten{\in}\mathrm X$, par la condition
\[
\begin{aligned}
&\handwritten{\varphi_{0}}(\mathrm U)=1\qquad\text{sss}\qquad \mathrm x\handwritten{\mathrm E}\mathrm U\,,\\
\text{i.e.}\qquad &\handwritten{\varphi}(\mathrm U)\ne\varnothing\quad \mathrel{\handwritten{\Leftarrow}\mkern-7mu=\mkern-7mu\handwritten{\Rightarrow}}\quad \mathrm x\handwritten{\in}\mathrm U\,.
\end{aligned}
\]
Soit $\mathrm F\handwritten{\in}\operatorname{Ob}\overset{\handwritten{\sim}}{\mathrm X_{\text{ét}}}$, alors l'image de $\mathrm F$ dans le faisceau final $\mathrm X$ est un ouvert; et comme
\[
\handwritten{\varphi}(\mathrm F)\mathrel{\mbox{-}\mkern-2mu\mbox{-}\mkern-3mu\handwritten{\to}}\handwritten{\varphi}(\mathrm U)
\]
est surj\makebox[0pt][l]{\raisebox{.95ex}{\handwritten{\scriptsize\unsure{''}}}}ectif, on voit que $\handwritten{\varphi}(\mathrm F)\ne\varnothing$ sss $\handwritten{\varphi}(\mathrm U)\ne\varnothing$ i.e. sss $\mathrm x\handwritten{\in}\mathrm U$. Soit alors $\mathrm C_{\handwritten{\varphi}}$ la catégorie des couples $(\mathrm X',\handwritten{\xi})$ où $\mathrm X'\handwritten{\in}\operatorname{Ob}\mathrm X_{\text{ét}}$, et $\handwritten{\xi\in\varphi}(\mathrm X')$. On a signalé dans \handwritten{7.9.1.} que $\mathrm C_{\handwritten{\varphi}}^{\handwritten{\unsure{o|0|\circ}}}$ est pseudo-filtrante connexe. De plus grâce à \handwritten{7.9.2.} et à ce qui précéde appliqué à $\mathrm X'$ au lieu de $\mathrm X$, si $(\mathrm X',\handwritten{\xi})\handwritten{\in}\operatorname{Ob}\mathrm X_{\text{ét.}}$, il lui est\unsure{:} associé un unique point $\mathrm x'\handwritten{\in}\mathrm X'$, tel que pour un morphisme étale $\mathrm X''\mathrel{\mbox{-}\mkern-2mu\mbox{-}\mkern-2mu\mbox{-}\mkern-3mu\handwritten{\to}}\mathrm X'$, $\handwritten{\xi}$ est dans l'image de $\handwritten{\varphi}(\mathrm X'')\mathrel{\mbox{-}\mkern-2mu\mbox{-}\mkern-3mu\handwritten{\to}}\handwritten{\varphi}(\mathrm X')$ sss $\mathrm x'$ est dans celle de $\mathrm X''$. On conclut de ceci que le\deleted{\unsure{s}} préschéma\deleted{\unsure{s}} $\underset{\handwritten{\longleftarrow}}{\lim}$ des $\mathrm X'$ suivant la catégorie $\mathrm C$ des $(\mathrm X',\handwritten{\xi})$ existe et est un localisé strict $\mathrm Z$ de $\mathrm X$, correspondant donc à un point géométrique $\handwritten{\xi}$ et un foncteur \deleted{\unsure{l|b}}fibre $\handwritten{{\unsure{\bar{\mathcal E}|\mathcal E|\bar{\varepsilon}|\varepsilon}}_{\xi}}$.
\end{minipage}
\end{lrbox}
\noindent\begin{tikzpicture}
\node[inner sep=0pt,outer sep=0pt] (body) {\begingroup\renewcommand{\ULthickness}{0pt}\deleted{\usebox{\coteFourTypescript}}\endgroup};
\draw[gray,thin] (body.south west) ++(.37\linewidth,.5em) -- ([xshift=-.31\linewidth,yshift=-2.6em]body.north east);
\end{tikzpicture}
% archive sheet number: 3
